%% nam-csound.tex — "Integrating Neural Amp Modeling into Csound"
%% NIME paper template (nimeart.cls), submission/review mode.
%% Converted from paper-draft.md on 2026-07-09.
%%
%% For camera-ready, change to: \documentclass[sigconf]{nimeart}
\documentclass[sigconf,anonymous,review]{nimeart}

%% This package avoids urls margin overflow
\usepackage{xurl}

%% \BibTeX command to typeset BibTeX logo in the docs
\AtBeginDocument{%
  \providecommand\BibTeX{{%
    Bib\TeX}}}

\setcopyright{cc}
\copyrightyear{2026}
\acmYear{2026}
\acmDOI{}

\acmConference[NIME '26]{International Conference on New Interfaces for Musical Expression}{June 23--26,
  2026}{London, UK}
%% This suppresses the ACM Reference Format printing.
\settopmatter{printacmref=false}
\acmISBN{}

\begin{document}

\title{Integrating Neural Amp Modeling into Csound}

%% Anonymous submission: the [anonymous,review] class options hide the
%% author block. Replace with real authors for camera-ready.
\author{Anonymous Author(s)}
\email{anonymous@example.com}
\affiliation{%
  \institution{Anonymous Institution}
  \city{Anonymous}
  \country{Anonymous}
}

\renewcommand{\shortauthors}{Anon.}

\begin{abstract}
  Neural Amp Modeler (NAM) is the premier open-source system for the
  black-box modeling of nonlinear audio hardware using neural networks. A
  community library of several hundred thousand pre-trained captures
  exists, but no way to use them inside Csound. This paper presents
  \texttt{NAMProcess}, a Csound opcode that performs real-time NAM
  inference. The opcode loads standard \texttt{.nam} model files, switches
  between models during performance without interrupting audio, adds no
  latency beyond the block boundary, and runs comfortably in real time on
  systems like the Raspberry Pi. This paper details the implementation of
  a thin wrapper around the NeuralAmpModelerCore C++ library and its
  integration as a native plugin opcode, a Cabbage plugin, a live
  dry-referenced volume compensation system, and ``tone selection'' in a
  Raspberry Pi-based guitar pedal.
\end{abstract}

\keywords{Neural Amp Modeling, Volume Compensation, Raspberry Pi, Cabbage
  Plugin, Csound Plugin Opcode}

\maketitle

\section{Introduction}

Over the past several years a new approach has become standard practice in
guitar audio: record a real device, then train a neural network to
reproduce its input--output behavior. Neural Amp Modeler
(NAM)~\cite{atkinson2025slimmable, nam} is the most widely used
open-source implementation of this approach. Its model format has become a
de facto interchange standard with over 350,000 accumulated on TONE3000
alone~\cite{tone3000} --- covering amplifiers, effect pedals, preamps, and
other outboard hardware, freely downloadable.

\texttt{NAMProcess} wraps the existing NeuralAmpModelerCore inference
library~\cite{nam} rather than reimplementing any DSP.
Section~\ref{sec:opcode} covers the opcode's design, including a loader
that never blocks the audio thread, correct handling of model state in
stereo, and measured latency. Section~\ref{sec:volume} presents a live
volume compensation system that references the player's dry signal, making
heterogeneous community captures usable side by side.
Sections~\ref{sec:plugin} and~\ref{sec:pedal} carry the engine into two
deployment targets, a drag-and-drop Cabbage plugin for DAW hosts and
``tone selection'' in a Raspberry Pi-based ``guitar pedal.''

\section{Background and Related Work}

Black-box neural modeling of guitar amplifiers was established by
Damsk\"agg et al., first against simulated tube-amplifier
data~\cite{damskagg2019tube} and then against real distortion circuits in
real time~\cite{damskagg2019distortion}, using a feedforward variant of
the WaveNet convolutional architecture. Wright et al.\ showed that
recurrent (LSTM) models reach comparable accuracy at lower
cost~\cite{wright2019rnn}, and consolidated the approach across amplifiers
and pedals~\cite{wright2020realtime}, training with an error-to-signal
ratio loss under a perceptual pre-emphasis
filter~\cite{wright2020perceptual}. Steinmetz and Reiss extended efficient
neural modeling to analog dynamic range
compression~\cite{steinmetz2022compressor}. NAM~\cite{atkinson2025slimmable,
nam} packages this family of techniques --- its model format supports
WaveNet and LSTM architectures --- together with a training pipeline whose
input is a fixed published reference signal reamped through the target
device.

Composers have extended music programming languages from within since the
MUSIC-N family itself --- Tenney's instrument building at Bell
Labs~\cite{tenney1969cme} sits at the head of the lineage that leads to
Csound --- and an opcode continues that practice with a learned model as
the unit generator.

Real-time neural inference has precedent in other computer music
environments. Chowdhury's RTNeural~\cite{chowdhury2021rtneural} provides a
general inference library designed for audio callbacks, and UGens built on
it exist for SuperCollider, with ports to Max/MSP and Pure
Data~\cite{pluta_rtneural}. The \texttt{neural\_tilde} external loads
\texttt{.nam} files directly in Max/MSP~\cite{neuraltilde}. Within Csound
the closest prior work is the component-level triode opcode of Fink and
Rabenstein~\cite{fink2011triode}, which models one circuit stage
physically rather than capturing a whole device from data. We present a
Csound~\cite{lazzarini2016csound} equivalent not present in the opcode
manual, the plugin-opcode collections, or community archives.

\section{The NAMProcess Opcode}
\label{sec:opcode}

\subsection{Interface}
\label{sec:interface}

\begin{verbatim}
aoutL, aoutR NAMProcess ainL, ainR,
      kProfileIdx, SLibDir
      [, kAutoVol, kDriveDb]
icount       NAMCount
\end{verbatim}

\texttt{SLibDir} names a directory scanned recursively for \texttt{.nam}
files at initialization. The sorted file list defines stable indices.
\texttt{kProfileIdx} selects the active model at k-rate. Writing a new
index triggers an asynchronous load, and the running model is replaced
only when the new one is ready. \texttt{NAMCount} returns the library size
so orchestra logic can wrap or clamp index changes. Because model choice
is an ordinary k-rate parameter, timbre selection is open to the same
rule-based and generative processes as any other Csound value ---
selection procedures in the sense Koenig practiced --- rather than
confined to manual preset browsing. \texttt{kAutoVol}
enables the volume compensation of Section~\ref{sec:volume} and defaults
to on. \texttt{kDriveDb} is a pre-model gain in dB whose loudness cost the
compensation absorbs, making it a pure saturation control. Existing
four-argument call sites compile unchanged. A minimal live instrument:

\begin{verbatim}
instr 1
  aL, aR  ins
  kIdx    chnget "model" ; MIDI, OSC, score
  aoL,aoR NAMProcess aL, aR, kIdx,
          "/home/user/models"
  outs aoL, aoR
endin
\end{verbatim}

Until the first model finishes loading, and whenever the directory
contains no models, the opcode passes audio through unchanged, so a
misconfigured path degrades to bypass rather than silence.

\subsection{Loader design}
\label{sec:loader}

All file and network-weight handling runs on a background thread owned by
the opcode, following standard real-time audio
discipline~\cite{bencina2011realtime}. The audio thread never allocates,
parses JSON, or blocks on I/O. Per k-cycle it takes one uncontended mutex
acquisition to copy a pointer to the active model pair. The loader thread
builds the requested model off the audio thread, selects the full-quality
submodel for slimmable models~\cite{atkinson2025slimmable}, and pre-warms
internal buffers before publishing the swap. A load failure retains the
previous model, so a corrupt file cannot interrupt a performance. Model
switching is consequently click-safe and non-blocking. The old model runs
until the exact block boundary at which the new pair is published.

When a model's training sample rate differs from the engine rate, the
loader wraps it in a Lanczos resampling container (the same mechanism the
reference plugin uses), so the network always runs at its trained rate.
Measured on a 48~kHz model under a 44.1~kHz engine, the harmonic spectrum
matches the matched-rate reference within 0.3~dB; without the wrapper the
model's learned filtering is stretched by the rate ratio and its whole
voicing shifts audibly.

\subsection{Model state and stereo processing}
\label{sec:stereo}

NAM models are stateful. WaveNet layers keep per-layer convolution history
in ring buffers, and LSTM models carry hidden state across process calls.
\texttt{NAMProcess} loads two independent instances of the same model
file, resets both off-thread, and publishes them atomically as a pair.
Each channel owns one instance. The cost is one extra load per switch and
double the model memory. It buys exact channel isolation. Processing both
channels through a shared instance instead convolves each block against
the other channel's immediate past, which measures as $-4.7$~dB crosstalk
on decorrelated stereo material.

\subsection{Registering an a-rate plugin opcode in Csound 6.18}
\label{sec:oentry}

One implementation detail is worth recording for other opcode developers.
Registered conventionally for an audio-rate generator, \texttt{OENTRY}
with \texttt{thread = 5}, the a-rate callback is never scheduled under
Csound 6.18's plugin registration path and the opcode emits silence with
no error. The working registration is \texttt{thread = 3}, computing the
full \texttt{ksmps} block inside the k-rate callback. Declaring
\texttt{outypes = "aa"} still allocates a-rate outputs.

\subsection{Latency}
\label{sec:latency}

Because the supported architectures are causal, the opcode should add no
delay beyond the \texttt{ksmps} block granularity that every Csound opcode
shares. Offline impulse renders confirm this. Response onset lies within
one sample of the input impulse. Total added latency is therefore one
k-cycle --- 128 samples, 2.7~ms at 48~kHz, in the pedal deployment of
Section~\ref{sec:pedal}.

\section{Live Volume Compensation}
\label{sec:volume}

The opcode continuously matches the wet output to the perceived loudness
of the player's dry signal. Load any capture --- clean, cranked,
mislabeled --- and the volume does not move.

Without this, browsing a community library is unusable live: fed the same
signal, 25 representative captures span about 10~dB, and their embedded
loudness metadata does not predict the ranking (the quietest-tagged
capture measures second-loudest). No pre-computed correction can close the
gap, because saturating models compress: a loudness measurement is only
valid near the input level at which it was taken, and one
compression-dominated capture holds its output within 1~dB across a 14~dB
input sweep, so no per-model gain equalizes it at all. The correction has
to be relative to what the player is actually doing, and it has to be
live.

Dry input and wet output are metered as ITU-R BS.1770 K-weighted mean
squares averaged with a 2.5~s time constant, and a slow output trim tracks
their ratio. Program dynamics move both meters together and leave the
ratio untouched, so the mechanism is not a compressor. The result is an
engineered perceptual invariance --- loudness held constant while spectrum
and saturation move --- the kind of controlled percept Tenney built pieces
around. Measured after
convergence, a $-12$~dB input step passes through as $-10.2$~dB of output,
the residual being the amp's own compression. The trim engages after one
second of active signal, slews at most 2~dB/s, holds within a 0.75~dB
deadband, freezes below $-55$~dBFS input, and is clamped to $\pm 18$~dB.
Converged trims are remembered per model file, so returning to a
previously played capture lands at level immediately ($+0.6$~dB measured,
inside the deadband). A load-time probe seeds the trim for unseen models;
the live loop owns the result.

Because the dry reference is tapped before the input stage,
\texttt{kDriveDb} pushes the model into saturation without a level change.
$+18$~dB of drive measured $+0.1$~dB of loudness and $-4.6$~dB of crest
factor. Across the 25-capture library on identical material, the best
static probe left a 9.9~dB spread (standard deviation 2.60~dB). The
dry-referenced loop leaves 3.1~dB (0.83~dB), with every capture within
$\pm 2$~dB of the dry signal, including the compression-dominated capture
that no input gain can equalize. The official NAM plugin's normalized
output mode applies a static gain from stored metadata or a fixed target.
To our knowledge no published NAM host performs live dry-referenced
compensation.

\section{A Drag-and-Drop Cabbage Plugin}
\label{sec:plugin}

We present the packaging of this opcode as a VST3/AU plugin through
Cabbage~\cite{walsh2011cabbage} not only to demonstrate our opcode's
integration into a plugin but to encourage an open-source platform for NAM
development in Csound. The plugin window is a single drop target. Cabbage
publishes a dropped file's path on its reserved
\texttt{LAST\_FILE\_DROPPED} channel, and a path-input variant of the
loader hot-swaps to the dropped model through the background mechanism of
Section~\ref{sec:loader}. The previous amp keeps running until the new one
is ready. Companion opcodes report the path of the actually-loaded model,
so the UI shows real state rather than the last request, and step through
the sorted \texttt{.nam} siblings of the active file for in-folder
browsing. The volume compensation of Section~\ref{sec:volume} defaults on,
with a toggle exposing raw model level for A/B comparison. The input
control is the loudness-compensated drive of Section~\ref{sec:interface},
so it shapes saturation rather than volume. A gate, three-band tone stack,
and DC blocker mirror the outboard chain of the reference NAM plugin.

The plugin is architected to degrade gracefully rather than go silent. A
plugin whose orchestra fails to compile emits silence yet still passes
host validation, so the base orchestra contains no NAM references at all.
The engine instrument is compiled at runtime with \texttt{compilestr}. If
the opcode library fails to load, audio passes through dry and the failure
is reported on the status display. The plugin also disables Csound's
default MIDI instrument assignment with \texttt{massign 0, 0}. The loader
registry is global, and a MIDI-spawned duplicate instrument instance would
otherwise free the loader out from under it.

\section{Tone Selection in a Raspberry Pi Guitar Pedal}
\label{sec:pedal}

The opcode's original deployment target is a Raspberry Pi 5-based guitar
pedal running Csound 6.18 at 48~kHz with \texttt{ksmps} 128, the audio
chain pinned to a dedicated core. A library of 25 captures acts as a bank
of different ``tones''. Model selection is mapped to MIDI foot controls
and OSC, arriving at the opcode as the k-rate index of
Section~\ref{sec:interface}. Because loading is asynchronous and
click-safe, the player steps between amps mid-phrase. Because compensation
runs inside the opcode, no volume ride accompanies a tone change, even for
brand-new additions to the library. Per-model trim memory keeps repeated
tone-stepping stable. Added latency is the one k-cycle of
Section~\ref{sec:latency} (2.7~ms), well inside a live monitoring budget.
CPU load with a full-quality WaveNet-family model has been informally
observed around 15\% of the audio core during performance.

\section{Conclusions}
\label{sec:conclusions}

\texttt{NAMProcess} gives Csound direct access to NAM's model format and
its community library. Any of several hundred thousand captures --- amps,
pedals, preamps, full chains --- becomes a drop-in nonlinear processor
addressable by a k-rate index, composable with everything else in the
orchestra. The same engine ships as a drag-and-drop Cabbage plugin and as
the tone section of a Raspberry Pi guitar pedal.

Several extensions are planned. Only one \texttt{NAMProcess} instance per
orchestra is currently supported, and lifting this is planned, as is an
n-channel variant following Csound's array conventions, per-model trim
memory persisted across sessions, and support for the plugin mechanisms of
Csound 7. Source code and measurement tooling are available at [link
withheld for review].

%% The Ethical Standards section is mandatory for all NIME submissions.
\section{Ethical Standards}

%% TODO: confirm/extend for camera-ready (funding, conflicts of interest).
This work did not involve human participants or animals. The community
model captures used for evaluation are freely downloadable; no potential
conflicts of interest are known to the authors.

%% Acknowledgments withheld for anonymous review; restore for camera-ready.
% \begin{acks}
% ...
% \end{acks}

\bibliographystyle{ACM-Reference-Format}
\bibliography{references}

\end{document}

%% ---------------------------------------------------------------------
%% Submission notes (carried over from paper-draft.md):
%%   ICSC 2026 deadline 2026-07-11 23:59 Rome time. llncs.cls template,
%%   4 pages max, 5th page references-only OK. Blind: AuthorA/AuthorB
%%   placeholders. Plugin name "NAM Drop", real repo link, acknowledgments
%%   (Dr. B co-authorship) at camera-ready only. Optional figure: plugin
%%   screenshot in a DAW (ICSC norm), crop name for blind. Verify Walsh
%%   cite form ("Ways Ahead", Cambridge Scholars 2013) before camera-ready.
%%   Static-calibration table (per-model input gains) removed 2026-07-09
%%   per author: live dry-referenced comp IS the method; static schemes
%%   appear only as compressed motivation. Pedal + plugin library = the
%%   25-capture set in cabbage-plugin/models/. Section 4 numbers all from
%%   volume-compensation-writeup.md (verified against shipped sources
%%   csound_nam_opcode.cpp / nam_loader.cpp / kweight.h, 2026-07-09).
%%   Full 9-approach experiment history lives in cabbage-plugin/archive/
%%   volume-compensation-history.md - repo material, not paper.
%% ---------------------------------------------------------------------
